\chapter{Examiner Defense Pack}
\label{app:examiner_defense_pack}

\noindent\textbf{Purpose.} This appendix consolidates the dissertation's key defenses against anticipated examiner challenges. It is intended for quick reference during viva preparation and provides a one-page-per-topic summary of how the dissertation addresses each anticipated question.

\section{Research Problem}
The central problem is that clinician trust in AI-assisted ASD-EEG screening is insufficient to support operational adoption, despite high classifier accuracy. The trust deficit is governance-driven rather than performance-driven (see Chapter~1, Section~\ref{sec:problem_statement}).

\section{Research Gap}
Existing research evaluates explainability, trust, and adoption independently. No prior study has integrated governance maturity as a measurable antecedent within a unified ASD-focused framework (see Chapter~2, Section~\ref{sec:governance_missing_variable}).

\section{Research Objectives}
Five objectives O1--O5 spanning B2C trust, technical validation, clinical adoption, governance maturity, and combined deployment verdict (see Chapter~1, Section~\ref{sec:research_aim}).

\section{Research Questions}
Five RQs aligned to objectives, evaluating whether governance maturity influences explainability, trust, and adoption (see Chapter~1, Section~\ref{sec:hypotheses}).

\section{Hypotheses}
Six DBA hypotheses H1--H6 covering the governance$\rightarrow$explainability$\rightarrow$trust$\rightarrow$adoption pathway (see Chapter~4, Section~\ref{sec:hypothesis_decision_summary}).

\section{Contributions}
Three primary contributions: governance-mediated adoption model (theory), RGAIG Maturity Framework (artifact), governance maturity as measurable construct (method) (see Chapter~5, Section~\ref{sec:executive_summary_contributions}).

\section{Theory Extension}
The dissertation extends TAM by positioning governance maturity as an upstream organizational antecedent (see Chapter~2, Section~\ref{sec:extending_tam_for_ai}; Chapter~5, Section~\ref{sec:theoretical_contribution_ch5}).

\section{Why DBA Not PhD CS}
The central focus is organizational adoption, not algorithmic innovation. ML/XAI/RAG serve as supporting infrastructure (see Chapter~1, Section~\ref{sec:why_dba_problem}; Chapter~3, Section~\ref{sec:separation_research_engineering}).

\section{Why PLS-SEM}
Selected for predictive theory testing with moderate sample sizes and mediation analysis support (see Chapter~3, Section~\ref{sec:why_pls_sem}).

\section{Why n=131}
Exceeded power analysis minimum and PLS-SEM guidance (see Chapter~3, Section~\ref{sec:why_n131}).

\section{Limitations}
Healthcare context specificity; intention-based adoption measurement; evolving governance perceptions; rapid GenAI advances (see Chapter~5, Section~\ref{sec:ch5_limitations}).

\section{Future Research}
Longitudinal designs, multi-country samples, real-world deployment, alternative AI modalities (see Chapter~5, Section~\ref{sec:ch5_future_research}).

\section{Novelty Statement}
First integrated empirical evaluation of governance maturity, explainability, trust, and adoption readiness within ASD-focused AI-assisted screening (see Chapter~1, Section~\ref{sec:doctoral_novelty_statement}; Chapter~5, Section~\ref{sec:what_makes_novel}).
